\documentclass[a4paper]{article}
\usepackage{graphicx}
\usepackage{amsmath,amssymb}
\usepackage{hyperref}
\usepackage{minted}
\usepackage{multirow}
\usepackage{float}
\input{/home/donkarlo/Dropbox/repo/nd_ai_project/src/nd_ai/language/formal/computer/programming/tex/composition/control_sequence/kind/tag/tag.tex}

\title{Development stages}
\date{}

\begin{document}
    \maketitle
    \url{https://en.wikipedia.org/wiki/Self-awareness#Human_development}
    \section{Levels}
        \begin{itemize}
            \item Level 0—\href{https://en.wikipedia.org/wiki/Confusion}{Confusion}: The person is unaware of any mirror reflection or the mirroring itself; they perceive a mirror image as an extension of their environment.
        \end{itemize}
        \begin{itemize}
            \item Level 1—Differentiation: The individual realizes the mirror is able to reflect things. They see that what is in the mirror is of a different nature from what is surrounding them. At this level they can differentiate between their own movement in the mirror and the movement of the surrounding environment.
        \end{itemize}
        \begin{itemize}
            \item Level 2—Situation: The individual can link the movements on the mirror to what is perceived within their own body.
        \end{itemize}
        \begin{itemize}
            \item Level 3—\href{https://en.wikipedia.org/wiki/Identification_(psychology)}{Identification}: An individual can now see that what's in the mirror is not another person but actually them.
        \end{itemize}
        \begin{itemize}
            \item Level 4—Permanence: The individual is able to identify the self in previous pictures looking different or younger. A "permanent self" is now experienced.
        \end{itemize}
        \begin{itemize}
            \item Level 5—\href{https://en.wikipedia.org/wiki/Self-consciousness}{Self-consciousness} or "meta" self-awareness: At this level, not only is the self seen from a first person view but it is realized that it is also seen from a third person's view. A person who develops self consciousness begins to understand they can be in the mind of others: for instance, how they are seen from a public standpoint.
        \end{itemize}

    % \bibliographystyle{plain}
    % \bibliography{/home/donkarlo/Dropbox/repo/shared_research/bibliography/bibliography.bib}
\end{document}